% ============================================================
%  Section 6: Tools -- XAMPP & phpMyAdmin
% ============================================================
\section{Tools for Database Development}

% ------------------------------------------------------------
\begin{frame}
  \frametitle{XAMPP -- Overview}

  \begin{block}{What Is XAMPP?}
    \textbf{XAMPP} is a free, open-source software bundle that installs a fully functional local web server and database environment in minutes -- no manual configuration required.
    \begin{itemize}
      \item Developed by \textbf{Apache Friends} since \textbf{2002}
      \item Available for \textbf{Windows}, \textbf{macOS}, and \textbf{Linux} (hence ``cross-platform'')
      \item Used in this course to host a local \textbf{MySQL\,/\,MariaDB} database for all exercises
    \end{itemize}
  \end{block}

  \begin{block}{The Acronym Explained}
    \begin{tabular}{@{}cl@{}}
      \textbf{X} & Cross-platform \\
      \textbf{A} & Apache HTTP Server \\
      \textbf{M} & MariaDB (originally MySQL) \\
      \textbf{P} & PHP \\
      \textbf{P} & Perl \\
    \end{tabular}
  \end{block}

  \begin{alertblock}{Why XAMPP?}
    Zero configuration needed -- every component works out of the box. Ideal for local development and course exercises where setting up a production server would be overkill.
  \end{alertblock}
\end{frame}

% ------------------------------------------------------------
\begin{frame}
  \frametitle{Installing XAMPP}

  \begin{block}{Download \& Install}
    \begin{enumerate}
      \item Download the installer from \textbf{apachefriends.org} (choose the version matching your OS)
      \item Run the installer -- \textit{accepting the default paths is strongly recommended}
      \item Default installation directories:
        \begin{itemize}
          \item Windows: \texttt{C:\textbackslash xampp}
          \item macOS: \texttt{/Applications/XAMPP}
        \end{itemize}
    \end{enumerate}
  \end{block}

  \begin{block}{After Installation: The Control Panel}
    The \textbf{XAMPP Control Panel} is your central dashboard. It lists all bundled services:
    \begin{itemize}
      \item Apache, \textbf{MySQL}, FileZilla, Mercury, Tomcat
    \end{itemize}
    To start the database: click \textbf{``Start''} next to \textit{MySQL}.\\
    The status indicator turns \textcolor{green!60!black}{\textbf{green}} when the service is running successfully.
  \end{block}

  \begin{examples}{Typical First-Time Pitfall}
    On Windows, port conflicts with Skype or IIS are common. If Apache or MySQL fail to start, check whether another process is already occupying port 80 (Apache) or port 3306 (MySQL).
  \end{examples}
\end{frame}

% ------------------------------------------------------------
\begin{frame}
  \frametitle{Starting and Accessing XAMPP}

  \begin{block}{Starting the Database Service}
    \begin{itemize}
      \item Open the XAMPP Control Panel and click \textbf{Start} next to \textit{MySQL}
      \item Default port: \textbf{3306} (the MySQL/MariaDB standard port)
      \item Default credentials:
        \begin{itemize}
          \item Username: \texttt{root}
          \item Password: \textit{(empty by default)}
        \end{itemize}
    \end{itemize}
  \end{block}

  \begin{alertblock}{Security Warning}
    An empty root password is acceptable \textit{only} on a local development machine with no external network access. \textbf{Never} deploy XAMPP with default credentials in a production or publicly accessible environment!
  \end{alertblock}

  \begin{block}{Access Options}
    \begin{itemize}
      \item \textbf{phpMyAdmin} (browser-based): navigate to \texttt{http://localhost/phpmyadmin}
      \item \textbf{MySQL command line}: available directly from the XAMPP shell
      \item \textbf{External SQL clients}: MySQL Workbench, DBeaver, DataGrip -- connect to \texttt{localhost:3306}
    \end{itemize}
    \textit{For our exercises: use phpMyAdmin or MySQL Workbench.}
  \end{block}
\end{frame}
